\documentclass[12pt]{article}
\usepackage[utf8]{inputenc}
\usepackage[T1]{fontenc}
\usepackage{hyperref}
\title{Watubela Language Resources for K-12 Curriculum Development}
\author{Compiled by Claude Sonnet 5 (claude-sonnet-5) for Andr\'as Kornai}
\date{}
\begin{document}
\maketitle

\section{Introduction}
Watubela (ISO 639-3 \texttt{wah}) is an Austronesian language of the Watubela Islands, East Seram Regency, Maluku, Indonesia. It is closely related to Geser and Bati. The figure of about 4,000 speakers dates from 1990. Verbix lists the dialects Sulmelang and Tamher Timur. Little educational material was found.

\section{Primary Educational Resources}
None were found.

\section{Linguistic and Reference Materials}
\begin{itemize}
\item Wikipedia article on the Watubela language: \url{https://en.wikipedia.org/wiki/Watubela_language}
\item Verbix language page, with dialect and lexical similarity data: \url{https://docs.verbix.com/Languages/Watubela}
\item Related language, Geser (Wikipedia): \url{https://en.wikipedia.org/wiki/Geser_language}
\end{itemize}

\section{Religious and Community Resources}
Joshua Project notes that Scripture and Christian materials are not yet available in Watubela.
\begin{itemize}
\item Joshua Project, Watubela, Wesi: \url{https://joshuaproject.net/people_groups/15891}
\end{itemize}

\section{Digital and Audio Resources}
None were found.

\section{Curriculum Development Context}
No school programme in Watubela was found. Indonesian is the likely language of instruction. Any curriculum would need new development with the community, possibly drawing on comparative material from Geser.
\end{document}